\item Una cuerda horizontal puede transmitir una potencia máxima $P_0$ (sin romperse) si por ella viaja una onda con amplitud $A$ y frecuencia angular $\omega$. Para aumentar esta potencia máxima, un estudiante dobla la cuerda y usa esta “cuerda doble” como medio. Determine la potencia máxima que se puede transmitir a lo largo de la “cuerda doble”, si supone que la tensión en las dos hebras juntas es la misma que la tensión original en la cuerda individual, y también que la frecuencia angular de la onda permanece constante.